\section{Related Work}%

% \ruijie{
As summarized in Table~\ref{tab:related_articulation}, we divide existing articulated reconstruction approaches into three types: optimization-based, learning-based, and generation-based methods.

\paragraph{Optimization-based articulated reconstruction.} 
An important line of work studies articulated objects through per-instance reconstruction from multiple observations. Early methods~\cite{liu2023paris,tseng2022cla,wei2022self} recover geometry and articulation by fitting a deformable NeRF~\cite{mildenhall2020nerf} to a specific object observed across multiple views or articulation states. More recent dynamic reconstruction methods~\cite{liu2025artgs,wu2025reartgs,ai2026aim,shen2025gaussianart,liu2025videoartgs,guo2026articulat3d} further exploit motion videos or continuous interaction to disentangle moving parts via Gaussian Splatting~\cite{kerbl20233d}. However, they are still optimization-intensive pipelines whose quality depends on how well the observations cover the scene.
Therefore, a generalizable end-to-end articulation estimation is what physical-world modeling really needs: not only for fast inference but also for a full understanding of physical cues.

\paragraph{Learning-based articulated reconstruction.} 
Recent learning-based methods usually predict articulated structure from a single image~\cite{chen2024urdformer,li2026monoart} or a single 3D observation~\cite{che2024op,Goyal_2025_ICCV,li2025particulate,gao2025meshart,song2025magicarticulate,li2025urdfanything}. 
For example, Particulate~\cite{li2025particulate} takes a single point cloud as input, and outputs the part segmentation of the point cloud and the kinematic parameter of those parts. 
Compared with per-instance optimization, these feed-forward methods are much more scalable and can perform inference directly and quickly. 
However, these methods also have limitations.
Although they curate diverse articulation datasets for training, their generalization remains unsatisfactory, even within the same category. 
% This limitation becomes more severe in open-world settings, where object geometry and kinematic layouts can differ substantially from those seen during training. 
We argue that this is because they use only a single static state as conditioning for predicting action-based articulation, which makes the task itself ambiguous: a single observation may correspond to different articulation patterns, even within the same category.
Therefore, we propose using multi-state observations as conditioning and designing an action-wise framework to resolve the ambiguity. 
% Furthermore, we prove that, with the help of VLMs, even a single state input can be extended to multiple states as conditioning for our model, outperforming state-of-the-art methods in a fair comparison.

\paragraph{Generation-based articulated reconstruction.} 
Another direction focuses on constructing articulated assets by retrieval, composition, or generation rather than directly inferring articulation from visual observations. 
These methods~\cite{liusingapo,lei2023nap,liu2024cage,chen2025freeart3d,chen2025ArtiLatent,wu2026urdfanything_plus,Su_2025_CVPR} assemble reusable parts, invoke articulation-aware priors, or combine generation with structural reasoning to produce simulation-ready articulated assets. 
For example, URDF-Anything+ takes an image and a mesh as inputs and leverages DiT~\cite{peebles2023scalable,chen2025autopartgen} to autoregressively generate articulated objects.
Recent multimodal frameworks~\cite{wang2025kinematify,qiu2025articulate,lu2025dreamart,wu2025dipo} further advance articulated asset generation using generative priors from VLMs~\cite{hurst2024gpt,gemini_image_pro}. 
Despite their impressive generalization in open-world scenes, they suffer from slow inference, LLM illusions, and difficulty handling complex topology.
% }



% \paragraph{Reasoning from observed state change.} Our work is motivated by the observation that articulation is most reliably revealed by how an object changes across states. Prior methods have leveraged multiple observations mainly in optimization-heavy reconstruction pipelines~\cite{liu2023paris,tseng2022cla,wei2022self}, where motion evidence is absorbed through iterative fitting rather than encoded by a generalizable feed-forward predictor. In contrast, we study articulation inference from two observed 3D states of the same object within a single forward pass. This setting lies between single-state prediction and full reconstruction: it preserves factual motion evidence from state change while avoiding per-instance test-time optimization. More importantly, our method does not merely treat the second state as additional input. Instead, it uses shared part queries to explain both states with a common part ordering, and refines query-to-point reasoning through mask-guided attention feedback. The shared queries establish cross-state part correspondence, while the mask-guided refinement localizes moving parts in each state before kinematic decoding. This distinction separates our method from single-state feed-forward baselines, generic mask-based part segmentation methods, and optimization-based reconstruction pipelines, and underlies the stronger robustness and cross-scene generalization observed in our experiments.


\label{sec:related}
